\chapter{Sputum Culture}

\section*{What Is This Test?}
Sputum culture is a microbiologic test in which expectorated lower respiratory tract secretions are cultured on selective and nonselective media to identify bacterial, mycobacterial, or fungal pathogens causing pneumonia, bronchitis, or other lower respiratory tract infections. The specimen is evaluated microscopically by Gram stain for suitability (presence of <10 squamous epithelial cells and >25 polymorphonuclear leukocytes per low-power field indicates an adequate specimen, free of significant oropharyngeal contamination). Sputum culture has a sensitivity of 50--80\% and specificity of 80--95\% for detecting the causative organism in community-acquired pneumonia (CAP), with yield varying by specimen quality, prior antibiotic exposure, and the specific pathogen. Common pathogens identified include \textit{Streptococcus pneumoniae}, \textit{Haemophilus influenzae}, \textit{Moraxella catarrhalis}, \textit{Staphylococcus aureus}, \textit{Pseudomonas aeruginosa}, and members of the Enterobacterales family. The test also provides antimicrobial susceptibility data (MICs by broth microdilution or gradient diffusion) to guide targeted therapy and detect resistance mechanisms including extended-spectrum beta-lactamases (ESBL), carbapenemases, and MRSA.

\section*{Alternative Names}
Respiratory Culture; Sputum C\&S (Culture and Sensitivity); Lower Respiratory Culture; Expectorated Sputum Culture; Tracheal Aspirate Culture; Endotracheal Aspirate Culture

\section*{Principle}
The specimen is first evaluated by Gram stain to assess specimen adequacy and provide a preliminary etiologic diagnosis (gram-positive cocci in pairs suggesting pneumococcus; gram-negative coccobacilli suggesting \textit{Haemophilus}; gram-negative rods suggesting Enterobacterales or \textit{Pseudomonas}). The specimen is then inoculated onto a panel of agar media: sheep blood agar (supports growth of most respiratory pathogens and detects hemolysis patterns), chocolate agar (supports growth of fastidious organisms including \textit{H. influenzae} and \textit{Neisseria} species), MacConkey agar (selective for gram-negative rods and differentiates lactose fermenters from non-fermenters), and mannitol salt agar or colistin-nalidixic acid (CNA) agar (selective for gram-positive organisms). Plates are incubated at 35--37\textdegree{}C in 5\% CO$_2$ for 18--48 hours. Potential pathogens are identified by colony morphology, Gram stain, biochemical tests, MALDI-TOF mass spectrometry, or automated identification systems (VITEK 2, BD Phoenix, MicroScan). Antimicrobial susceptibility testing is performed by broth microdilution, disk diffusion (Kirby-Bauer), or gradient diffusion (Etest) according to CLSI or EUCAST breakpoints. Semi-quantitative reporting (rare, few, moderate, many or 1+ to 4+) correlates with the density of the pathogen in the specimen. Organisms present in quantities below the threshold for clinical significance or representing normal oropharyngeal flora are reported as "normal respiratory flora" or "mixed upper respiratory flora."

\section*{History}
The use of sputum examination for the diagnosis of tuberculosis was championed by Robert Koch in the 1880s. The Gram stain was developed by Hans Christian Gram in 1884 and quickly applied to sputum specimens for the diagnosis of pneumonia. The utility of sputum Gram stain and culture in community-acquired pneumonia was systematically evaluated by John Bartlett in the 1970s, establishing the foundational quality-assessment criteria still used today (<10 squamous epithelial cells/LPF, >25 WBC/LPF). The protected specimen brush (PSB) and bronchoalveolar lavage (BAL) were developed in the 1980s to improve diagnostic yield and avoid oropharyngeal contamination, especially in ventilated patients. Automated blood culture systems (BACTEC, BacT/ALERT) were adapted for lower respiratory specimens. MALDI-TOF mass spectrometry for bacterial identification, introduced in the 2010s, reduced identification time from 24--48 hours to minutes. Despite the ascendancy of molecular methods (multiplex PCR panels), sputum culture remains essential for antimicrobial susceptibility testing, detection of emerging resistance patterns, and recovery of organisms for epidemiologic typing.

\section*{Why This Test Is Ordered}
\begin{itemize}
  \item Diagnosis of the causative bacterial pathogen in patients with community-acquired pneumonia (CAP) requiring hospitalization, to allow targeted antimicrobial therapy and de-escalation from broad-spectrum empiric regimens
  \item Evaluation of healthcare-associated pneumonia (HCAP), hospital-acquired pneumonia (HAP), or ventilator-associated pneumonia (VAP) where multidrug-resistant pathogens (MRSA, \textit{P. aeruginosa}, \textit{Acinetobacter baumannii}, ESBL-producing Enterobacterales) are more prevalent
  \item Investigation of suspected bacterial exacerbation of chronic obstructive pulmonary disease (COPD) or bronchiectasis with increased sputum volume, purulence, and dyspnea
  \item Diagnosis of pulmonary tuberculosis by mycobacterial culture (sputum AFB culture) in patients with compatible clinical and radiographic findings
  \item Detection of fungal pathogens (\textit{Aspergillus} species, \textit{Histoplasma capsulatum}, \textit{Coccidioides immitis}, \textit{Blastomyces dermatitidis}) in immunocompromised patients with pneumonia not responding to antibiotics
  \item Identification of the pathogen and antimicrobial susceptibility in patients with lung abscess or empyema
  \item Antibiotic stewardship: targeted therapy based on identified pathogen and susceptibility results to reduce unnecessary broad-spectrum antibiotic exposure and development of resistance
\end{itemize}

\section*{Primary vs Secondary Test}
Sputum Gram stain and culture are primary diagnostic tests for hospitalized patients with community-acquired pneumonia (in accordance with IDSA/ATS guidelines recommendation for obtaining pretreatment sputum cultures in patients with severe CAP, those with risk factors for resistant pathogens, and those not responding to empiric therapy). Blood cultures are obtained concurrently. In non-severe outpatient CAP, sputum culture is generally not required. Multiplex PCR panels and urinary antigen testing complement but do not replace sputum culture, as culture provides antimicrobial susceptibility data.

\section*{How This Test Is Performed}
Specimen collection: The patient should rinse the mouth with water (not mouthwash) to reduce oral contamination, take a deep breath, and cough forcefully from deep within the chest to bring up sputum (not saliva). The expectorated specimen is collected directly into a sterile, leak-proof container. At least 1 mL of purulent material should be collected. Early morning specimens (first cough of the day) often have the highest yield. For patients unable to expectorate spontaneously, sputum induction with hypertonic saline (3\%) via nebulizer may be performed. In intubated patients, endotracheal aspirate or bronchoalveolar lavage (BAL) specimens are collected. Specimens should be transported to the laboratory at room temperature within 2 hours of collection, or refrigerated at 2--8\textdegree{}C if transport delay exceeds 2 hours. Anaerobic transport media is not required unless anaerobic infection is specifically suspected (lung abscess, empyema).

\section*{What Preparation Is Needed}
The patient should rinse the mouth with water immediately before specimen collection to minimize oropharyngeal contamination. The patient should not eat, drink, use mouthwash, or brush teeth for 1 hour before collection. For sputum induction, fasting for 2 hours before the procedure is recommended to reduce risk of aspiration. Antibiotics should ideally be withheld until after specimen collection, as even a single dose can reduce culture yield by 30--50\%.

\section*{Contraindications and Precautions}
Sputum induction with hypertonic saline is contraindicated in patients with severe asthma (risk of bronchospasm), untreated pneumothorax, hemoptysis > 50 mL, or recent thoracic surgery. Spontaneous expectoration has no absolute contraindications. Sputum specimens that consist primarily of saliva (more than 10 squamous epithelial cells per LPF on Gram stain) should be rejected as they represent oropharyngeal contamination, not lower respiratory tract secretions. Up to 30--50\% of specimens may be rejected for this reason. Previous antibiotic therapy significantly reduces culture sensitivity, particularly for fastidious organisms like \textit{S. pneumoniae}. The sputum culture cannot reliably diagnose anaerobic lung infections, as the specimen inevitably passes through the oropharynx which harbors abundant anaerobes.

\section*{Risks}
Spontaneous sputum expectoration carries no significant risk. Sputum induction may cause coughing, bronchospasm, or, rarely, hemoptysis. There is no risk from the laboratory testing itself. Healthcare workers handling specimens from patients with suspected tuberculosis must use Biosafety Level 2 practices with BSL-3 precautions for aerosol-generating procedures.

\section*{What Happens After the Test}
Preliminary Gram stain results are available within 1--2 hours of specimen receipt. Cultures are examined at 18--24 hours and final results, including identification and antimicrobial susceptibility, are available at 48--72 hours. Preliminary results may be communicated earlier if a dominant pathogen is identified. Empiric broad-spectrum antibiotics should be narrowed (de-escalated) to targeted therapy based on culture and susceptibility results. Negative culture at 48 hours in a patient who has not received antibiotics argues against a bacterial etiology and should prompt consideration of alternative diagnoses or discontinuation of antibiotics. Mycobacterial cultures are incubated for 6--8 weeks.

\section*{Understanding Results}

\subsection*{Below Normal}
\begin{itemize}
  \item \textbf{Normal respiratory flora} or \textbf{Mixed upper respiratory flora}: Indicates growth of organisms normally present in the oropharynx (\textit{viridans group streptococci}, \textit{Neisseria} species, coagulase-negative staphylococci, \textit{Corynebacterium} species, nonpathogenic \textit{Neisseria}) without a predominant pathogen. This result may represent: absence of bacterial lower respiratory tract infection; viral or non-infectious cause of respiratory symptoms; or a false-negative culture due to prior antibiotic therapy, inadequate specimen, or fastidious organism not recovered on routine media
  \item \textbf{No growth}: Indicates no bacterial growth on any media after 48 hours. Most commonly seen in patients pretreated with antibiotics or with viral pneumonia
\end{itemize}

\subsection*{Above Normal}
\begin{itemize}
  \item \textbf{Predominant pathogen identified with semi-quantitation (moderate to many, 3+ to 4+):} Indicates a causative bacterial pathogen, particularly when correlated with a compatible Gram stain result. Common patterns:
    \begin{itemize}
      \item \textit{Streptococcus pneumoniae} (gram-positive lancet-shaped diplococci): Most common cause of CAP. Blood culture should be obtained; urinary antigen testing can confirm
      \item \textit{Haemophilus influenzae} (gram-negative coccobacilli): Second most common cause of CAP in COPD patients. Beta-lactamase testing guides amoxicillin-clavulanate versus amoxicillin use
      \item \textit{Moraxella catarrhalis} (gram-negative diplococci): Common in COPD exacerbations. Typically beta-lactamase positive
      \item \textit{Staphylococcus aureus}: Post-influenza pneumonia, cavitary pneumonia, intravenous drug users. MRSA status critical for therapy selection
      \item \textit{Pseudomonas aeruginosa}: Hospital-acquired pneumonia, neutropenic patients, bronchiectasis, cystic fibrosis
      \item \textit{Klebsiella pneumoniae} or other Enterobacterales: Alcoholism, aspiration, healthcare-associated pneumonia
    \end{itemize}
  \item \textbf{Mixed growth with potential pathogen:} When a potential pathogen is isolated in mixed culture (1+ to 2+, or rarely present), clinical correlation is required. Significance is more likely if the same organism grows in blood culture
\end{itemize}

\section*{Normal Results}
\textbf{Normal respiratory flora} or \textbf{Mixed upper respiratory flora} or \textbf{No pathogenic bacteria isolated}. No growth on mycobacterial or fungal selective media.

\section*{Follow-up Tests}
\begin{itemize}
  \item \textbf{Blood cultures (2 sets):} Should be obtained before antibiotic therapy in all hospitalized patients with pneumonia; bacteremia confirms the causative pathogen identified in sputum
  \item \textit{\textbf{S. pneumoniae}} \textbf{urinary antigen:} Rapid immunochromatographic test with sensitivity 70--80\% and specificity >95\% for pneumococcal pneumonia, useful when sputum is unavailable or contaminated
  \item \textit{\textbf{Legionella pneumophila}} \textbf{urinary antigen:} When Legionella is suspected clinically (severe CAP, hyponatremia, GI symptoms, elevated LFTs, recent travel)
  \item \textbf{Mycobacterial culture (AFB):} When risk factors for tuberculosis, cavitary disease on imaging, or AFB seen on smear
  \item \textbf{Fungal culture and serum fungal markers (beta-D-glucan, galactomannan):} In immunocompromised patients with pneumonia not responding to antibiotics
  \item \textbf{Multiplex respiratory pathogen panel (PCR):} When viral or atypical bacterial etiology is suspected and sputum culture is negative
  \item \textbf{Chest CT:} When sputum culture is inconclusive and parenchymal disease characterization is needed
\end{itemize}

\section*{References}
\begin{enumerate}
  \item Metlay JP, Waterer GW, Long AC, et al. Diagnosis and treatment of adults with community-acquired pneumonia. An official clinical practice guideline of the ATS and IDSA. Am J Respir Crit Care Med. 2019;200(7):e45-e67.
  \item Bartlett JG. Diagnostic tests for agents of community-acquired pneumonia. Clin Infect Dis. 2011;52(Suppl 4):S296-S304.
  \item Mandell LA, Wunderink RG, Anzueto A, et al. IDSA/ATS consensus guidelines on the management of community-acquired pneumonia in adults. Clin Infect Dis. 2007;44(Suppl 2):S27-S72.
  \item Murray PR, Witebsky FG. The clinician and the microbiology laboratory. In: Mandell, Douglas, and Bennett's Principles and Practice of Infectious Diseases. 9th ed. Elsevier; 2019.
  \item CLSI. Principles and Procedures for Detection of Bacteria in Clinical Specimens. M40-A2. Clinical and Laboratory Standards Institute; 2014.
\end{enumerate}

\clearpage
\section*{Regulatory Status and Authorization Date}
Regulatory status applies to the specific assay, instrument, manufacturer, and intended use rather than to the test name alone. No single approval date can be assigned to this test category. For a commercial assay, verify the current product in the relevant regulator's database: the U.S. Food and Drug Administration (FDA) clearance, approval, or authorization; India's Central Drugs Standard Control Organisation (CDSCO) licence or permission; the European Union In Vitro Diagnostic Regulation (EU IVDR) CE certificate issued through the applicable conformity-assessment route; or the United Kingdom Medicines and Healthcare products Regulatory Agency (MHRA) registration and UKCA/CE status. Laboratory-developed tests may not have an individual FDA, CDSCO, EU IVDR, or MHRA approval date.
\clearpage
